% Options for packages loaded elsewhere
\PassOptionsToPackage{unicode}{hyperref}
\PassOptionsToPackage{hyphens}{url}
%
\documentclass[
  ignorenonframetext,
  aspectratio=169,
  russian,
]{beamer}
\usepackage{pgfpages}
\setbeamertemplate{caption}[numbered]
\setbeamertemplate{caption label separator}{: }
\setbeamercolor{caption name}{fg=normal text.fg}
\beamertemplatenavigationsymbolshorizontal
% Prevent slide breaks in the middle of a paragraph
\widowpenalties 1 10000
\raggedbottom
\setbeamertemplate{part page}{
  \centering
  \begin{beamercolorbox}[sep=16pt,center]{part title}
    \usebeamerfont{part title}\insertpart\par
  \end{beamercolorbox}
}
\setbeamertemplate{section page}{
  \centering
  \begin{beamercolorbox}[sep=12pt,center]{section title}
    \usebeamerfont{section title}\insertsection\par
  \end{beamercolorbox}
}
\setbeamertemplate{subsection page}{
  \centering
  \begin{beamercolorbox}[sep=8pt,center]{subsection title}
    \usebeamerfont{subsection title}\insertsubsection\par
  \end{beamercolorbox}
}
\AtBeginPart{
  \frame{\partpage}
}
\AtBeginSection{
  \ifbibliography
  \else
    \frame{\sectionpage}
  \fi
}
\AtBeginSubsection{
  \frame{\subsectionpage}
}

\usepackage{amsmath,amssymb}
\usepackage{iftex}
\ifPDFTeX
  \usepackage[T1]{fontenc}
  \usepackage[utf8]{inputenc}
  \usepackage{textcomp} % provide euro and other symbols
\else % if luatex or xetex
  \usepackage{unicode-math}
  \defaultfontfeatures{Scale=MatchLowercase}
  \defaultfontfeatures[\rmfamily]{Ligatures=TeX,Scale=1}
\fi
\usepackage{lmodern}
\usetheme[]{metropolis}
\ifPDFTeX\else  
    % xetex/luatex font selection
\fi
% Use upquote if available, for straight quotes in verbatim environments
\IfFileExists{upquote.sty}{\usepackage{upquote}}{}
\IfFileExists{microtype.sty}{% use microtype if available
  \usepackage[]{microtype}
  \UseMicrotypeSet[protrusion]{basicmath} % disable protrusion for tt fonts
}{}
\usepackage{xcolor}
\newif\ifbibliography
\setlength{\emergencystretch}{3em} % prevent overfull lines
\setcounter{secnumdepth}{-\maxdimen} % remove section numbering


\providecommand{\tightlist}{%
  \setlength{\itemsep}{0pt}\setlength{\parskip}{0pt}}\usepackage{longtable,booktabs,array}
\usepackage{calc} % for calculating minipage widths
\usepackage{caption}
% Make caption package work with longtable
\makeatletter
\def\fnum@table{\tablename~\thetable}
\makeatother
\usepackage{graphicx}
\makeatletter
\newsavebox\pandoc@box
\newcommand*\pandocbounded[1]{% scales image to fit in text height/width
  \sbox\pandoc@box{#1}%
  \Gscale@div\@tempa{\textheight}{\dimexpr\ht\pandoc@box+\dp\pandoc@box\relax}%
  \Gscale@div\@tempb{\linewidth}{\wd\pandoc@box}%
  \ifdim\@tempb\p@<\@tempa\p@\let\@tempa\@tempb\fi% select the smaller of both
  \ifdim\@tempa\p@<\p@\scalebox{\@tempa}{\usebox\pandoc@box}%
  \else\usebox{\pandoc@box}%
  \fi%
}
% Set default figure placement to htbp
\def\fps@figure{htbp}
\makeatother

\IfFileExists{plex-otf.sty}{
  %% Full TeXlive
  % \usepackage[%
  %   % math,
  %   RM={Scale=0.94},SS={Scale=0.94},SScon={Scale=0.94},TT={Scale=MatchLowercase,FakeStretch=0.9},DefaultFeatures={Ligatures=Common}
  % ]{plex-otf}
}{
  %% TinyTeX
  \usepackage{libertine}
}

%%% Load theme
% https://deic.uab.cat/~iblanes/beamer_gallery/
\IfFileExists{beamerthemegotham.sty}{
  %% Full TeXlive
  \usetheme{gotham}
  \gothamset{
    numbering=totalpagenumber,
    parttocframe default=off,
    sectiontocframe default=off,
    subsectiontocframe default=off,
  }
}{
  %% TinyTeX
  \usetheme{Madrid}
}
\metroset{progressbar=frametitle,sectionpage=progressbar,numbering=fraction}
\makeatletter
\@ifpackageloaded{caption}{}{\usepackage{caption}}
\AtBeginDocument{%
\ifdefined\contentsname
  \renewcommand*\contentsname{Содержание}
\else
  \newcommand\contentsname{Содержание}
\fi
\ifdefined\listfigurename
  \renewcommand*\listfigurename{Список иллюстраций}
\else
  \newcommand\listfigurename{Список иллюстраций}
\fi
\ifdefined\listtablename
  \renewcommand*\listtablename{Список таблиц}
\else
  \newcommand\listtablename{Список таблиц}
\fi
\ifdefined\figurename
  \renewcommand*\figurename{Рисунок}
\else
  \newcommand\figurename{Рисунок}
\fi
\ifdefined\tablename
  \renewcommand*\tablename{Таблица}
\else
  \newcommand\tablename{Таблица}
\fi
}
\@ifpackageloaded{float}{}{\usepackage{float}}
\floatstyle{ruled}
\@ifundefined{c@chapter}{\newfloat{codelisting}{h}{lop}}{\newfloat{codelisting}{h}{lop}[chapter]}
\floatname{codelisting}{Список}
\newcommand*\listoflistings{\listof{codelisting}{Листинги}}
\makeatother
\makeatletter
\makeatother
\makeatletter
\@ifpackageloaded{caption}{}{\usepackage{caption}}
\@ifpackageloaded{subcaption}{}{\usepackage{subcaption}}
\makeatother

\ifLuaTeX
\usepackage[bidi=basic]{babel}
\else
\usepackage[bidi=default]{babel}
\fi
\babelprovide[main,import]{russian}
\babelprovide[import]{english}
% get rid of language-specific shorthands (see #6817):
\let\LanguageShortHands\languageshorthands
\def\languageshorthands#1{}
\usepackage{csquotes}
\usepackage{bookmark}

\IfFileExists{xurl.sty}{\usepackage{xurl}}{} % add URL line breaks if available
\urlstyle{same} % disable monospaced font for URLs
\hypersetup{
  pdftitle={Лабораторная работа №1},
  pdfauthor={Савенкова Татьяна},
  pdflang={ru-RU},
  hidelinks,
  pdfcreator={LaTeX via pandoc}}


\title{Лабораторная работа №1}
\subtitle{Операционные системы}
\author{Савенкова Татьяна}
\date{Invalid Date}
\institute{Российский университет дружбы народов, Москва, Россия}

\begin{document}
\frame{\titlepage}


\section{Информация}\label{ux438ux43dux444ux43eux440ux43cux430ux446ux438ux44f}

\begin{frame}{Докладчик}
\phantomsection\label{ux434ux43eux43aux43bux430ux434ux447ux438ux43a}
\begin{columns}[c]
\begin{column}{0.7\linewidth}
\begin{itemize}[<+->]
\tightlist
\item
  Савенкова Татьяна Александровна
\item
  Студент НКАбд-05-25
\item
  я Таня
\item
  Российский университет дружбы народов
\item
  \href{mailto:1032253537@pfur.ru}{\nolinkurl{1032253537@pfur.ru}}
\end{itemize}
\end{column}

\begin{column}{0.3\linewidth}
\pandocbounded{\includegraphics[keepaspectratio]{./image/savenkova.png}}
\end{column}
\end{columns}
\end{frame}

\section{Цель
работы}\label{ux446ux435ux43bux44c-ux440ux430ux431ux43eux442ux44b}

\begin{frame}{Цель работы}
Получение навыков правильной работы с репозиториями git.
\end{frame}

\section{Задание}\label{ux437ux430ux434ux430ux43dux438ux435}

\begin{frame}{Задание}
\begin{itemize}[<+->]
\tightlist
\item
  Выполнить работу для тестового репозитория.
\item
  Преобразовать рабочий репозиторий в репозиторий с git-flow и
  conventional commits.
\end{itemize}
\end{frame}

\section{Теоретическое
введение}\label{ux442ux435ux43eux440ux435ux442ux438ux447ux435ux441ux43aux43eux435-ux432ux432ux435ux434ux435ux43dux438ux435}

\begin{frame}[fragile]{Теоретическое введение}
Общая информация

\begin{verbatim}
Gitflow Workflow опубликована и популяризована Винсентом Дриссеном.
Gitflow Workflow предполагает выстраивание строгой модели ветвления с учётом выпуска проекта.
Данная модель отлично подходит для организации рабочего процесса на основе релизов.
Работа по модели Gitflow включает создание отдельной ветки для исправлений ошибок в рабочей среде.
Последовательность действий при работе по модели Gitflow:
    Из ветки master создаётся ветка develop.
    Из ветки develop создаётся ветка release.
    Из ветки develop создаются ветки feature.
    Когда работа над веткой feature завершена, она сливается с веткой develop.
    Когда работа над веткой релиза release завершена, она сливается в ветки develop и master.
    Если в master обнаружена проблема, из master создаётся ветка hotfix.
    Когда работа над веткой исправления hotfix завершена, она сливается в ветки develop и master.
\end{verbatim}

Процесс работы с Gitflow

\begin{verbatim}
Основные ветки (master) и ветки разработки (develop)
    Для фиксации истории проекта в рамках этого процесса вместо одной ветки master используются две ветки. В ветке master хранится официальная история релиза, а ветка develop предназначена для объединения всех функций. Кроме того, для удобства рекомендуется присваивать всем коммитам в ветке master номер версии.

    При использовании библиотеки расширений git-flow нужно инициализировать структуру в существующем репозитории:

    git flow init

    Для github параметр Version tag prefix следует установить в v.

    После этого проверьте, на какой ветке Вы находитесь:

    git branch

Функциональные ветки (feature)
    Под каждую новую функцию должна быть отведена собственная ветка, которую можно отправлять в центральный репозиторий для создания резервной копии или совместной работы команды. Ветки feature создаются не на основе master, а на основе develop. Когда работа над функцией завершается, соответствующая ветка сливается обратно с веткой develop. Функции не следует отправлять напрямую в ветку master.
    Как правило, ветки feature создаются на основе последней ветки develop.

    Создание функциональной ветки

        Создадим новую функциональную ветку:

        git flow feature start feature_branch

        Далее работаем как обычно.

    Окончание работы с функциональной веткой

        По завершении работы над функцией следует объединить ветку feature_branch с develop:

        git flow feature finish feature_branch

Ветки выпуска (release)
    Когда в ветке develop оказывается достаточно функций для выпуска, из ветки develop создаётся ветка release. Создание этой ветки запускает следующий цикл выпуска, и с этого момента новые функции добавить больше нельзя — допускается лишь отладка, создание документации и решение других задач. Когда подготовка релиза завершается, ветка release сливается с master и ей присваивается номер версии. После нужно выполнить слияние с веткой develop, в которой с момента создания ветки релиза могли возникнуть изменения.
    Благодаря тому, что для подготовки выпусков используется специальная ветка, одна команда может дорабатывать текущий выпуск, в то время как другая команда продолжает работу над функциями для следующего.

    Создать новую ветку release можно с помощью следующей команды:

    git flow release start 1.0.0

    Для завершения работы на ветке release используются следующие команды:

    git flow release finish 1.0.0
    
    
\end{verbatim}

Ветки исправления (hotfix) Ветки поддержки или ветки hotfix используются
для быстрого внесения исправлений в рабочие релизы. Они создаются от
ветки master. Это единственная ветка, которая должна быть создана
непосредственно от master. Как только исправление завершено, ветку
следует объединить с master и develop. Ветка master должна быть помечена
обновлённым номером версии. Наличие специальной ветки для исправления
ошибок позволяет команде решать проблемы, не прерывая остальную часть
рабочего процесса и не ожидая следующего цикла релиза.

\begin{verbatim}
    Ветку hotfix можно создать с помощью следующих команд:

    git flow hotfix start hotfix_branch

    По завершении работы ветка hotfix объединяется с master и develop:

    git flow hotfix finish hotfix_branch
    
\end{verbatim}
\end{frame}

\section{Выполнение лабораторной
работы}\label{ux432ux44bux43fux43eux43bux43dux435ux43dux438ux435-ux43bux430ux431ux43eux440ux430ux442ux43eux440ux43dux43eux439-ux440ux430ux431ux43eux442ux44b}

\begin{frame}{Выполнение лабораторной работы}
Устанавливаю nodejs, пакетный менеджер для него pnpm и
gitflow.(рис.~\ref{fig-001}).

\begin{figure}

\centering{

\includegraphics[width=0.7\linewidth,height=\textheight,keepaspectratio]{image/1.png}

}

\caption{\label{fig-001}Установка необходимого ПО}

\end{figure}%

Устаналиваю через pnpm commitizen и standard-changelog.
(рис.~\ref{fig-002})

\begin{figure}

\centering{

\includegraphics[width=0.7\linewidth,height=\textheight,keepaspectratio]{image/2.png}

}

\caption{\label{fig-002}Установка через pnpm}

\end{figure}%

Создаю новый репозиторий и делаю там первый коммит. (рис.~\ref{fig-003})

\begin{figure}

\centering{

\includegraphics[width=0.7\linewidth,height=\textheight,keepaspectratio]{image/3.png}

}

\caption{\label{fig-003}Создание репозитория}

\end{figure}%

Инициализирую и конфигурирую общепринятые коммиты в созданной директории
через редактирование package.json. (рис.~\ref{fig-004})

\begin{figure}

\centering{

\includegraphics[width=0.7\linewidth,height=\textheight,keepaspectratio]{image/4.png}

}

\caption{\label{fig-004}Конфигурация package.json}

\end{figure}%

Делаю снимок изменений, создаю коммит и отправляю на удаленный
репозиторий. (рис~\ref{fig-005})

\begin{figure}

\centering{

\includegraphics[width=0.7\linewidth,height=\textheight,keepaspectratio]{image/5.png}

}

\caption{\label{fig-005}Отправление коммита на GitHub}

\end{figure}%

Инициализирую в репозитории git flow и создаю 1 релиз в только что
созданной ветке develop. (рис.~\ref{fig-006})

\begin{figure}

\centering{

\includegraphics[width=0.7\linewidth,height=\textheight,keepaspectratio]{image/6.png}

}

\caption{\label{fig-006}Использование git-flow}

\end{figure}%

Создаю список изменений через standard changelog, заканчиваю релиз и
выгружаю на удаленный репозиторий изменения. (рис.~\ref{fig-007})

\begin{figure}

\centering{

\includegraphics[width=0.7\linewidth,height=\textheight,keepaspectratio]{image/7.png}

}

\caption{\label{fig-007}Загрузка релиза на github}

\end{figure}%

Инициализирую ветку feature для работы над новой функциональностью,
готовлю релиз и загружаю на github. ({[}рис.@fig-008{]})

\begin{figure}

\centering{

\includegraphics[width=0.7\linewidth,height=\textheight,keepaspectratio]{image/8.png}

}

\caption{\label{fig-008}Работа с новой веткой и загрузка релиза}

\end{figure}%
\end{frame}

\section{Выводы}\label{ux432ux44bux432ux43eux434ux44b}

\begin{frame}{Выводы}
В ходе выполнения лабораторный работы я получила навыки правильной
работы с репозиториями git

:::
\end{frame}




\end{document}
